\subsection{Retrospective real-EEG check of the properly predictive comparator}
\label{sec:v18-fair-predictive}
The original class-centred comparator does not account for
between-source-person class-mean shifts
(Section~\ref{sec:v17-mean-shift-decomposition}, Eq.~\eqref{eq:v17-exact-nll-shift}).
We therefore executed a \emph{new secondary, post-outcome} comparison
of the correctly defined $C_{\mathrm{pred}}$ selector on public
Shin2017A EEG using the original, independent previously executed
preprocessing and classifier implementation. This additional
analysis is neither a new dataset nor a renewal of the earlier
confirmatory study. The Shin2017A outcomes had been seen before
the new score was specified.

For each of 29 outer held participants, the entire source-only
tangent reference and within-class state were reconstructed using
the other 28 participants. The source-participant-grouped inner
validation used only source-person trial labels; for each held
\emph{source} validation fold, the class means appearing in
Eq.~\eqref{eq:v16-conditional-prediction} were fitted exclusively
on the corresponding inner \emph{training} source groups.
The candidate set was that fold's originally source-feasible
27-point shrinkage grid, and the per-person proper-predictive
Gaussian log scores were averaged with equal source-person
weight. All 29 selected shrinkage strengths were serialised to
a machine-readable file and locked with SHA-256
\texttt{fac20e7f1ec1e43f3fa01af8ce88e44869e80f624d9a10b6123a57a6687f9b5c}
before re-opening the immutable original ZIP and reading the
outer target AUC outcomes. The original class-centred comparator
and source-macro policy were unchanged; the original Shin
score archive retained its previously verified SHA-256.

\begin{table}[tb]
\centering
\small
\caption{Shin2017A real EEG secondary comparison (\(N=29\);
1740 held-person test trials). The first and third rows were
part of the original frozen comparisons; the proper-predictive
row was newly run after the original target results were known.
It is \emph{descriptive/post-outcome} and does not enter
the original prospective Holm family. The proper-predictive
strength was 0.2 in all 29 folds.}
\label{tab:v18-real-predictive}
\begin{tabular}{lcc}
\toprule
Source selection objective & Mean AUC & Evidence tier\\
\midrule
Source participant macro-AUC & 0.785900 & Original frozen\\
Proper Gaussian class-conditional predictive log score &
 \textbf{0.786207} & New post-outcome\\
Validation-class-centred covariance score & 0.763678 & Original frozen\\
\midrule
Source-macro minus proper predictive & $-0.000307$ &
 Descriptive only\\
Class-centred minus proper predictive & $-0.022529$ &
 Descriptive only\\
\bottomrule
\end{tabular}
\end{table}

The newly implemented proper-predictive objective selected
\(\lambda=0.2\) in \emph{all 29} outer folds. It differs from the
original class-centred covariance choice in 29/29 folds, yet its
person-macro target AUC of \(0.786207\) was almost identical
to the source-macro policy's \(0.785900\). Proper predictive
scoring exceeded the original class-centred covariance
selection by \(0.022529\) person-macro AUC. In this
single previously exposed EEG dataset, the large comparison
against the class-centred score therefore does
\emph{not} translate into a discrimination-selector advantage
against genuine predictive likelihood.

The agreement between the two better-performing methods
should not be misrepresented as a general statistical
equivalence or proven non-inferiority. The new result
contains no new hypothesis test or adjusted $p$-value,
and equal source-population folds remain strongly dependent.
All 29 proper-predictive choices being 0.2 demonstrates
\emph{selection stability on highly overlapping training pools},
not successful participant-specific adaptation.
This new comparison narrows the paper's conclusion to
the \emph{measurement target itself}: a held-source
class-centred covariance log score can omit the
class-mean-shift penalty, changing the selected
regularisation regime and biasing the interpretation
of EEG decoding method comparisons.
